\begin{table}[htbp]
\centering
\caption{Summary Statistics}
\label{tab:sumstat}
\begin{tabular}{lccccc}
\hline\hline
 & Mean & SD & Min & Max & $N$ \\
\hline
\multicolumn{6}{l}{\textit{Panel A: Patent Variables}} \\[3pt]
Patent grants (county-year)                   & 118.5 & 456.8 & 0 & 11,865 & 7,824 \\
Bartik instrument                             & 0.7 & 0.1 & 0 & 1 & 7,824 \\
\\[6pt]
\multicolumn{6}{l}{\textit{Panel B: Employment Outcomes ($t+1$)}} \\[3pt]
Employment (thousands)                        & 107.7 & 228.7 & 2 & 3,854 & 7,824 \\
New hires (thousands, annual)                 & 69.5 & 146.3 & 0 & 2,845 & 7,824 \\
Monthly earnings (\$)                         & 3,338.3 & 716.1 & 1,924 & 9,141 & 7,824 \\
\\[6pt]
\multicolumn{6}{l}{\textit{Panel C: Sector Employment ($t+1$, thousands)}} \\[3pt]
Exposed sectors (Mfg + Info + Prof/Sci)       & 22.8 & 54.3 & 0 & 1,054 & 7,824 \\
Local services (Retail + Accom + Other)       & 29.4 & 59.4 & 0 & 1,059 & 7,824 \\
\hline\hline
\multicolumn{6}{p{0.95\textwidth}}{\begin{flushleft}\small
Notes: Sample consists of 876 US counties with at least 20 patent applications during the 2001--2003 share-construction window, observed annually over 2004--2012. Patent grants are counted by first-inventor county using PatentsView geocoded locations linked to USPTO Patent Examination Research data. Employment outcomes are from the Census Bureau Quarterly Workforce Indicators (QWI), aggregated from quarterly to annual frequency (employment: mean of four quarters; hires: sum). The Bartik instrument is the county-level shift-share instrument constructed from pre-determined (2001--2003) county art-unit patent application shares interacted with leave-one-out examiner leniency shocks.
\end{flushleft}}
\end{tabular}
\end{table}
